% Part: normal-modal-logic
% Chapter: axioms-systems
% Section: normal-logics

\documentclass[../../../include/open-logic-section]{subfiles}

\begin{document}

\olfileid[hi]{nml}{prf}{nor}

\olsection{सामान्य मोडल तर्कशास्त्र}

मोडल !!{formula}ों के प्रत्येक समुच्चय को अभिगृहीतों के किसी समुच्चय से
व्युत्पाद्य !!{formula}ों के रूप में आसानी से अभिलक्षित नहीं किया जा सकता। हम
चाहते हैं कि मोडल तर्कशास्त्र सुव्यवस्थित हों। सबसे पहले, शास्त्रीय
प्रतिज्ञप्तिक तर्कशास्त्र में जो कुछ हम !!{derive} वह अब भी
!!{derivable} होना चाहिए, यह ध्यान रखते हुए कि सूत्रों में अब
\iftag{prvBox}{$\Box$\iftag{prvDiamond}{
    और~}{}}{}\iftag{prvDiamond}{$\Diamond$}{} भी आ सकते हैं। इसके लिए हम
अपेक्षा करते हैं कि मोडल तर्कशास्त्र सभी सर्वसत्यतात्मक दृष्टांतों को समाहित
करे और मोडस पोनेन्स के अधीन संवृत हो।

\begin{defn}
  \emph{मोडल तर्कशास्त्र} मोडल !!{formula}ों का ऐसा समुच्चय~$\Sigma$ है जो
  \begin{enumerate}
  \item सभी सर्वसत्यताओं को समाहित करता है, और
  \item प्रतिस्थापन के अधीन संवृत है; अर्थात् यदि $!A \in \Sigma$ और
    $!D_1$, \dots, $!D_n$ !!{formula} हैं, तो
    \[
    \SSubst{!A}{\subst{!D_1}{p_1}, \dots, \subst{!D_n}{p_n}} \in \Sigma,
    \]
    \item \emph{मोडस पोनेन्स} के अधीन संवृत है; अर्थात् यदि $!A$ और
      $!A \lif !B \in \Sigma$, तो $!B \in \Sigma$।
  \end{enumerate}
\end{defn}

मोडल तर्कशास्त्र के संबंधात्मक अर्थविज्ञान का उपयोग करने के लिए हमें यह भी
अपेक्षा करनी होती है कि सभी मोडल प्रतिरूपों में मान्य सभी !!{formula} समाहित
हों। यह अपेक्षा तब पूरी हो जाती है जब \Ax{K}\iftag{prvDiamond}{ और~\Dual{}}{}
के सभी दृष्टांत !!{derivable} हों और जब भी कोई !!{formula}~$!A$
!!{derivable} हो, $\Box !A$ भी हो। इन प्रतिबंधों को संतुष्ट करने वाला मोडल
तर्कशास्त्र \emph{सामान्य} कहलाता है। (निस्संदेह गैर-सामान्य मोडल
तर्कशास्त्र भी हैं, पर सामान्य संबंधात्मक प्रतिरूप उनके लिए पर्याप्त नहीं।)

\begin{defn}
  कोई मोडल तर्कशास्त्र $\Sigma$ \emph{सामान्य} है यदि उसमें
  \begin{align*}
    \tag{\Ax{K}} & \Box(p \lif q) \lif (\Box p \lif \Box q),
    \iftag{prvDiamond}{\\
      \tag{\Dual} & \Diamond p \liff \lnot\Box\lnot p}{}
  \end{align*}
  हों और वह \emph{आवश्यकीकरण} के अधीन संवृत हो; अर्थात् यदि
  $!A \in \Sigma$, तो $\Box !A \in \Sigma$।
\end{defn}

ध्यान दें कि सर्वसत्यतात्मक निहितार्थ \emph{किसी जगत पर सत्यता} को संरक्षित
रखने जितना ``सूक्ष्म'' है, जबकि नियम \Nec{} केवल \emph{किसी प्रतिरूप में
सत्यता} को संरक्षित रखता है (और इसलिए किसी फ़्रेम या फ़्रेमों के वर्ग में
मान्यता को भी)।

\begin{prop}\ollabel{prop:rk}
  प्रत्येक सामान्य मोडल तर्कशास्त्र नियम \RK{} के अधीन संवृत है:
  \begin{prooftree}
    \AxiomC{$!A_1 \lif (!A_2 \lif \cdots (!A_{n-1} \lif !A_n)\cdots)$}
    \RightLabel{\RK}
    \UnaryInfC{$\Box!A_1 \lif (\Box!A_2 \lif \cdots (\Box!A_{n-1}
      \lif \Box!A_n)\cdots).$}
  \end{prooftree}
\end{prop}

\begin{proof}
  $n$ पर आगमन से: यदि $n = 1$, तो नियम केवल \Nec{} है और प्रत्येक सामान्य
  मोडल तर्कशास्त्र \Nec{} के अधीन संवृत है।

  अब मान लें परिणाम $n-1$ के लिए मान्य है; हम इसे~$n$ के लिए दिखाते हैं।

  मान लें
  \begin{align*}
  & {}!A_1 \lif (!A_2 \lif \cdots (!A_{n-1} \lif !A_n)\cdots) \in \Sigma
  \intertext{आगमन-परिकल्पना से हमारे पास है:}
  & \Box!A_1 \lif (\Box!A_2 \lif \cdots \Box(!A_{n-1} \lif !A_n)\cdots)
  \in \Sigma
  \intertext{चूँकि $\Sigma$ सामान्य मोडल तर्कशास्त्र है, उसमें $\Ax{K}$ के
    सभी दृष्टांत हैं, विशेषतः:}
  & \Box(!A_{n-1} \lif !A_n) \lif (\Box!A_{n-1} \lif \Box!A_n) \in \Sigma
  \intertext{मोडस पोनेन्स और उपयुक्त सर्वसत्यतात्मक दृष्टांतों से हमें मिलता है:}
  & \Box!A_1 \lif (\Box!A_2 \lif \cdots (\Box!A_{n-1}
  \lif \Box!A_n)\cdots) \in \Sigma. 
  \end{align*}
\end{proof}

\begin{prop}\ollabel{prop:notDiamondBot}
  प्रत्येक सामान्य मोडल तर्कशास्त्र $\Sigma$ में~$\lnot\Diamond\lfalse$ है।
\end{prop}

\begin{prob}
  \olref[nml][prf][nor]{prop:notDiamondBot} सिद्ध कीजिए।
\end{prob}

\begin{prop}
  मान लें $!A_1$, \dots, $!A_n$ !!{formula} हैं। तब $!A_1$, \dots, $!A_n$
  के सभी दृष्टांतों को समाहित करने वाला एक सबसे छोटा मोडल तर्कशास्त्र
  $\Sigma$ है।
\end{prop}

\begin{proof}
  $!A_1$, \dots, $!A_n$ दिए होने पर $\Sigma$ को उन सभी सामान्य मोडल
  तर्कशास्त्रों का सर्वनिष्ठ परिभाषित करें जिनमें $!A_1$, \dots, $!A_n$ के
  सभी दृष्टांत हैं। सर्वनिष्ठ अरिक्त है, क्योंकि सभी !!{formula}ों का समुच्चय
  $\Frm[L]$ ऐसा मोडल तर्कशास्त्र है।
\end{proof}

\begin{defn}
$!A_1$, \dots, $!A_n$ को समाहित करने वाला सबसे छोटा सामान्य मोडल तर्कशास्त्र
\emph{मोडल तंत्र} कहलाता है और $\Log{K} !A_1 \dots !A_n$ से निरूपित होता
है। सबसे छोटा सामान्य मोडल तर्कशास्त्र~\Log{K} से निरूपित होता है।
\end{defn}

\end{document}
